\chapter{Stabilisation Policy and Its Limits}\label{ch:b1:stabilisation-policy}

Between 1992 and 2014, 63 economies went through 153 recessions, counting
a recession as a year in which output fell. In April of the year before,
forecasters expected output to fall in 5 of them; by October of the
recession year itself, in 118. % ledger: b1:an2018
The last two chapters showed how policy can move output in a model. This
one asks what stands between the model and a policy that works: late
information, slow decisions, effects that arrive later still, and
uncertainty about the model itself.

\begin{openingchart}{Change in the general-government balance (points of GDP)
against the change in real GDP growth (points), 2007 to 2009, 27 European
Union countries, with the fitted line.}{Eurostat (gov\_10dd\_edpt1,
nama\_10\_gdp)}
\begin{tikzpicture}
\begin{axis}[omscatter, xlabel={Change in GDP growth, 2007 to 2009 (points)},
  ylabel={Change in budget balance (points of GDP)}, xmin=-28, xmax=-2, ymin=-16, ymax=2,
  width=0.62\linewidth, height=0.45\linewidth]
  \addplot[only marks, mark=*, mark size=1.5pt, omDef] table[x=dgr, y=dbal, col sep=comma] {figdata/bachelor-1/26-stabilisation-policy/out/budget.csv};
  \addplot[only marks, mark=none, nodes near coords, point meta=explicit symbolic, nodes near coords style={font=\scriptsize, anchor=west}]
    table[x=dgr, y=dbal, meta=geo, col sep=comma] {figdata/bachelor-1/26-stabilisation-policy/out/budget_labels.csv};
  \addplot[omThm, thick, sharp plot, mark=none] table[x=x, y=fit, col sep=comma] {figdata/bachelor-1/26-stabilisation-policy/out/budget_fit.csv};
\end{axis}
\end{tikzpicture}
\end{openingchart}

From 2007 to 2009 the budget balance fell in 26 of the 27 countries. The
fitted line rises by only $0.122$ points of GDP per point of growth lost, with
a correlation of $0.235$: the size of the fall in growth explains little of
the size of the fall in the balance. % figdata: budget_stats
Part of each fall came by itself, as tax receipts fell and benefits rose;
part came from decisions and from revenue tied to particular sectors ---
Ireland's balance fell by $14.2$ points and Spain's by $13.1$, for falls in growth
of $10.4$ and $7.3$ points. % figdata: budget

\begin{recall}
Monetary policy, rules and \omterm{def:b1:stabilisation-policy:rule}{discretion}, the Taylor rule, fiscal policy and
public debt, automatic stabilisers and the output gap were met at school.
IS--LM is \cref{ch:b1:is-lm}; AD--AS and the Phillips curve are
\cref{ch:b1:aggregate-supply}.
\end{recall}

\section{Delays}\label{sec:b1:stabilisation-policy:lags}

A shock happens; it shows up in the data a quarter or two later, and in
revised data later still; a decision is taken; the measure is put in
place; and its effect on spending builds over several quarters.

\begin{definition}[Policy lag]\label{def:b1:stabilisation-policy:lags}
\index{policy lag}\index{inside lag}\index{outside lag}
A \emph{policy lag} is the time between a shock and the effect of the policy
that responds to it. The \emph{inside lag} runs from the shock to the action
(recognition, decision, implementation); the \emph{outside lag} from the
action to its effect on the economy.
\end{definition}

Monetary policy has a short \omterm{def:b1:stabilisation-policy:lags}{inside lag} --- a committee meets every few
weeks --- and a long \omterm{def:b1:stabilisation-policy:lags}{outside lag}; fiscal policy usually the reverse.

\begin{proposition}[A lagged response can amplify a cycle]\label{prop:b1:stabilisation-policy:cycle}
Let the output gap without policy be a cycle $x_t = A\sin(2\pi t/T)$, and let policy add
$-\theta x_{t-k}$, a response of size $\theta > 0$ to the gap of $k$ periods earlier. Then the gap
with policy is a cycle of the same period whose variance is
\[
1 + \theta^2 - 2\theta\cos(2\pi k/T)
\]
times the variance without policy. Policy raises the variance when $\cos(2\pi k/T) < \theta/2$.
\end{proposition}

\begin{proof}
$x_t - \theta x_{t-k} = A[\sin\phi_t - \theta\sin(\phi_t - \psi)]$ with $\phi_t = 2\pi t/T$ and $\psi = 2\pi k/T$. Expanding,
this is $A[(1 - \theta\cos\psi)\sin\phi_t + \theta\sin\psi\cos\phi_t]$, a sine of amplitude
$A\sqrt{(1 - \theta\cos\psi)^2 + \theta^2\sin^2\psi} = A\sqrt{1 + \theta^2 - 2\theta\cos\psi}$. The variance of a sine is half the
squared amplitude, so the ratio is $1 + \theta^2 - 2\theta\cos\psi$, which exceeds 1 when $\cos\psi < \theta/2$.
\end{proof}

\begin{omfigure}
\begin{tikzpicture}
\begin{axis}[omrecord, grid=none, xlabel={Quarter}, ylabel={Output gap (\%)}, xmin=0, xmax=16, ymin=-2, ymax=2,
  width=0.62\linewidth, height=0.38\linewidth, samples=97, domain=0:16, xtick={0,4,...,16}]
  \addplot[black, thick] {sin(deg(2*pi*x/8))};
  \addplot[omThm, thick] {sin(deg(2*pi*x/8)) - 0.5*sin(deg(2*pi*(x-4)/8))};
  \addplot[omDef, thick, dashed] {sin(deg(2*pi*x/8)) - 0.5*sin(deg(2*pi*(x-1)/8))};
  \node[font=\scriptsize, anchor=south] at (axis cs:10,1.5) {lag 4: amplitude $1.5$};
  \node[font=\scriptsize, anchor=west] at (axis cs:1.6,-1.7) {dashed, lag 1: amplitude $0.74$};
\end{axis}
\end{tikzpicture}

{\small A cycle of eight quarters (black) and the same cycle with a policy
response of half the gap ($\theta = 0.5$). With a lag of one quarter (dashed blue) the
variance falls to $0.543$ of its value; with a lag of four quarters, half the
cycle (red), policy pushes in the wrong direction every quarter and the
variance rises to $2.25$ times.}
\end{omfigure}

\begin{model}{Stabilisation with lags}
\index{stabilisation policy}
\assumes The economy of \cref{ch:b1:is-lm,ch:b1:aggregate-supply}; policy
responds to the gap of $k$ periods earlier, by an amount $\theta$.
\predicts A countercyclical policy can raise the volatility of output
when the lag is long relative to the cycle, or when the response is too
strong for the persistence of the shocks.
\tested Forecasts made in April of the year before missed 148 of 153
recessions (An, Jalles and Loungani 2018); % ledger: b1:an2018
policy recommendations computed from real-time data differ considerably
from those computed with revised data (Orphanides 2001). % ledger: b1:orphanides2001
The lags and errors the model assumes are real.
\breaks When policy responds to forecasts rather than to past gaps; when
its effect on expectations is immediate, as with announcements of future
policy; and when the cycle is not a regular wave of fixed length.
\end{model}

\section{Uncertainty about the model}\label{sec:b1:stabilisation-policy:uncertainty}

Even with the gap known, the effect of policy is not. Let the gap be $x + mp$,
where $p$ is policy and $m$ its effect, a random variable with mean $\mu$ and
variance $\sigma^2$, and let the loss be the expected squared gap.

\begin{proposition}[Act less under uncertainty]\label{prop:b1:stabilisation-policy:brainard}
$E[(x + mp)^2] = (x + \mu p)^2 + \sigma^2p^2$ is minimised at $p^* = -\mu x/(\mu^2 + \sigma^2)$, smaller in absolute value
than the response $-x/\mu$ that closes the gap on average.
\end{proposition}

\begin{proof}
The derivative is $2\mu(x + \mu p) + 2\sigma^2p$, zero at $p^*$; since $\mu^2/(\mu^2 + \sigma^2) < 1$, $|p^*| < |x/\mu|$.
\end{proof}

The more uncertain the effect, the more of the gap it pays to leave open:
a full response adds its own risk. With a gap of $-4$, $\mu = 1$ and $\sigma^2 = 0.5$, the full
response $p = 4$ leaves an expected loss of $8$; the optimal $p^* = 2.667$ leaves $5.333$;
doing nothing, $16$. % computed: code/bachelor-1/26-stabilisation-policy/exercises/solutions.py

\section{The record of fine-tuning}\label{sec:b1:stabilisation-policy:finetuning}

\begin{record}[Fine-tuning]
In the 1960s it was widely held that policy could keep output close to its
potential quarter by quarter, by adjusting spending, taxes and money as
forecasts required. The following decade combined high inflation with
recessions (\cref{ch:b1:aggregate-supply}). Estimates made in real time,
the kind policy relies on, can differ considerably from those computed
once the data are revised, and so can the policy they recommend
(Orphanides 2001). % ledger: b1:orphanides2001
The claim that policy could fine-tune output was refuted: the lags, the
errors in real-time data and the response of expectations were all larger
than assumed.
\end{record}

\section{Automatic stabilisers}\label{sec:b1:stabilisation-policy:automatic}

Income taxes fall when incomes fall, and benefits rise when people lose
their jobs. Both soften a fall in income without a decision, and so
without an \omterm{def:b1:stabilisation-policy:lags}{inside lag}. Their size depends on how progressive the tax is
and how generous the benefits; the opening chart shows the balance falling
in nearly every European country in 2008--09, partly through them.

Brenland's household survey shows how much (weekend problem): a fall of
$10\%$ in the earnings of every working household lowers \omterm{def:b1:national-accounts-in-depth:gni}{disposable income} by
only $68.8\%$ of the fall in market income. Taxes absorb $24.7\%$ of it and
transfers $6.5\%$. % computed: code/bachelor-1/26-stabilisation-policy/exercises/solutions.py

\section{Rules and discretion}\label{sec:b1:stabilisation-policy:rules}

\begin{definition}[Policy rule and discretion]\label{def:b1:stabilisation-policy:rule}
\index{policy rule}\index{discretion}
A \emph{policy rule} fixes in advance how a policy instrument responds to
the state of the economy. Policy is made with \emph{discretion} when the
instrument is set case by case, with no commitment to a rule.
\end{definition}

\omterm{def:b1:stabilisation-policy:rule}{Discretion} can respond to shocks
a rule did not foresee; a rule can be known in advance, which helps the
public form its expectations, and is harder to bend. The best-known
example is the Taylor rule (Taylor 1993), % ledger: b1:taylor1993
\[
i = \pi + 0.5(y - y^*) + 0.5(\pi - 2) + 2,
\]
where $i$ is the policy rate, $\pi$ inflation and $y - y^*$ the output gap, all in per
cent. It raises the real rate when inflation is above $2\%$ or output above
potential. A rule needs the gap, and the gap is known only with error; the
rule inherits the error (\cref{exo:b1:stabilisation-policy:realtime}).

A further argument for rules --- that a policy-maker free to choose each
period has an incentive to surprise, which the public anticipates --- is
time inconsistency, a subject of the second university year.

\section{How multipliers are measured}\label{sec:b1:stabilisation-policy:multipliers}

The IS--LM multiplier of \cref{ch:b1:is-lm} depends on slopes; the
evidence uses three designs. \emph{Narrative shocks} use changes in
spending decided for reasons unrelated to the state of the economy, such
as wars. \emph{Regional cross-sections} compare regions that received
different amounts of the same national spending. \emph{Structural models}
estimate the slopes and compute the multiplier.

\begin{definition}[Local multiplier]\label{def:b1:stabilisation-policy:local-multiplier}
\index{local multiplier}\index{relative multiplier}
A \emph{local} (or \emph{relative}) \emph{multiplier} is the effect on a region's output,
relative to other regions, of spending in that region, relative to other
regions. National monetary policy and national taxes are common to all
regions and drop out of the comparison.
\end{definition}

\begin{debate}[How large are fiscal multipliers?]
With quarterly American data covering several large wars and deep
recessions, Ramey and Zubairy (2018) find spending multipliers below 1,
whether or not the economy has slack; near the \omterm{def:b1:is-lm:liquidity-trap}{zero lower bound} the
results are mixed, a few specifications giving 1.5. % ledger: b1:ramey2018
With regime-switching models on post-war data, Auerbach and Gorodnichenko
(2012) find multipliers much larger in recessions than in expansions.
% ledger: b1:auerbach2012
Using regional differences in military build-ups, Nakamura and Steinsson
(2014) estimate a \omterm{def:b1:stabilisation-policy:local-multiplier}{local multiplier} of about 1.5. % ledger: b1:nakamura2014
Each design has its weakness: wars are few and differ from peacetime
spending; regime-switching results depend on how a recession is dated;
a \omterm{def:b1:stabilisation-policy:local-multiplier}{local multiplier} is not the national one, since it differences out the
monetary and tax response. What the evidence has refuted is a large
multiplier, fixed at all times, of the kind the Keynesian cross gives.
\end{debate}

\begin{omchart}{Consumer prices, United States, 12-month change, monthly,
2019--2024, with the three largest relief packages marked.}{Bureau of
Labor Statistics, CPI-U all items; package dates and sizes from Barnichon,
Oliveira and Shapiro (2021)}
\begin{tikzpicture}
\begin{axis}[omtime, xlabel={Year}, ylabel={CPI, 12-month change (\%)}, xmin=2019, xmax=2025, ymin=0, ymax=10,
  width=0.66\linewidth, height=0.38\linewidth]
  \addplot+[thick, no marks, omDef] table[x=t, y=yoy, col sep=comma] {figdata/bachelor-1/26-stabilisation-policy/out/cpi.csv};
  \draw[gray, dashed] (axis cs:2020.21,0) -- (axis cs:2020.21,8.4);
  \draw[gray, dashed] (axis cs:2020.96,0) -- (axis cs:2020.96,8.4);
  \draw[gray, dashed] (axis cs:2021.21,0) -- (axis cs:2021.21,8.4);
  \node[font=\scriptsize, anchor=south, align=center] at (axis cs:2020.21,8.4) {2.2};
  \node[font=\scriptsize, anchor=south, align=center] at (axis cs:2020.9,8.4) {0.9};
  \node[font=\scriptsize, anchor=south, align=center] at (axis cs:2021.27,8.4) {1.9};
\end{axis}
\end{tikzpicture}
\end{omchart}

\begin{record}[2020--23]
The United States enacted relief of 2.2 trillion dollars in March 2020, 0.9
trillion more in December 2020 and 1.9 trillion in March 2021, the
dashed lines of the chart (trillions of dollars above them).
% ledger: b1:frbsf-packages
Consumer-price inflation rose to $9.1\%$ over the twelve months to June
2022, then fell to $2.9\%$ by December 2024. % figdata: cpi
The first transfers supported incomes when output fell; the later ones came
when the gap was closing.
\end{record}

\begin{debate}[What caused the inflation of 2021--23?]
Estimates of the stimulus's effect range widely, and measure different
things. One put the
effect of the third package alone at about $0.3$ points of inflation a year
through 2022 (Barnichon, Oliveira and Shapiro 2021); % ledger: b1:frbsf-packages
another, comparing the United States with other advanced economies, found
that all pandemic fiscal support may have raised American inflation by
about 3 points by the end of 2021, with considerable uncertainty (Jordà,
Liu, Nechio and Rivera-Reyes 2022). % ledger: b1:jlnr2022
Decompositions of the inflation itself also differ. Sorting price changes
by whether prices and quantities moved together, one study attributes
about half of the run-up to supply factors and about a third to demand,
the rest being ambiguous (Shapiro 2022); % ledger: b1:shapiro2022
another finds that most of the surge came from shocks to prices given
wages --- commodity prices pushed up by strong demand, and shortages in
particular sectors --- rather than from tight labour markets, whose effect
was smaller but more persistent (Blanchard and Bernanke 2023). % ledger: b1:bb2023
The estimates use different designs, periods and definitions of demand;
none settles the share of each cause.
\end{debate}

\section*{Exercises}\label{sec:b1:stabilisation-policy:exercises}

\begin{exercise}[$\star$]\label{exo:b1:stabilisation-policy:three}
Output starts falling in January. The statistics office reports the fall
in April; the parliament votes a spending package in September; the
spending starts in January of the next year and has most of its effect by
the summer. Name and date the lags.
\end{exercise}

\begin{exercise}[$\star$]\label{exo:b1:stabilisation-policy:read}
In the opening chart's table, Germany's balance fell by $3.4$ points for a fall
in growth of $8.5$ points, and Ireland's by $14.2$ for a fall of $10.4$.
% figdata: budget
Compare the two, and say what besides automatic stabilisers can explain
the difference.
\end{exercise}

\begin{exercise}[$\star$]\label{exo:b1:stabilisation-policy:taylor}
With the Taylor rule of \cref{sec:b1:stabilisation-policy:rules}, find the
policy rate when inflation is $4\%$ and the output gap $-1\%$.
\end{exercise}

\begin{exercise}[$\star$]\label{exo:b1:stabilisation-policy:sign}
In a recession with no discretionary change in policy, predict the sign of
the change in the budget balance, and say why.
\end{exercise}

\begin{exercise}[$\star\star$]\label{exo:b1:stabilisation-policy:cycle-ex}
With the cycle of \cref{prop:b1:stabilisation-policy:cycle}, $T = 8$ and $\theta = 0.5$,
compute the variance ratio for lags of 1, 2 and 4 quarters.
\end{exercise}

\begin{exercise}[$\star\star$]\label{exo:b1:stabilisation-policy:uncertain}
The gap is $-4$, the effect of policy has mean 1 and variance $0.5$. Find the
optimal response and the expected loss of the optimal response, of the full
response and of no response.
\end{exercise}

\begin{exercise}[$\star\star$]\label{exo:b1:stabilisation-policy:slope}
The fitted line of the opening chart has slope $0.122$ and correlation $0.235$.
% figdata: budget_stats
What would the slope be if the balance moved only through stabilisers that
offset a fixed share of the fall in output? Why is the measured slope a
poor measure of their size?
\end{exercise}

\begin{exercise}[$\star\star$]\label{exo:b1:stabilisation-policy:realtime}
A central bank follows the Taylor rule, but its real-time estimate of the
gap is $-3\%$ when the true gap is $-1\%$. By how much is its rate too low, and
what does the Phillips curve predict for inflation if the error lasts?
\end{exercise}

\begin{exercise}[$\star\star$]\label{exo:b1:stabilisation-policy:local}
A study finds a \omterm{def:b1:stabilisation-policy:local-multiplier}{local multiplier} of 1.5. Give two reasons why the national
multiplier of the same spending could be smaller, and one why it could be
larger.
\end{exercise}

\begin{exercise}[$\star\star\star$]\label{exo:b1:stabilisation-policy:ar}
The gap follows $y_t = \rho y_{t-1} + \varepsilon_t + p_t$, with $\varepsilon_t$ independent with variance $\sigma^2$
and $0 < \rho < 1$. Policy sets $p_t = -\theta y_{t-1}$. Find the variance of $y$, and show that
policy raises it if $\theta > 2\rho$.
\end{exercise}

\begin{exercise}[$\star\star\star$]\label{exo:b1:stabilisation-policy:rule-disc}
Under \omterm{def:b1:stabilisation-policy:rule}{discretion} the loss from each period is $4$, the cost of an inflation
bias the public expects. A rule removes the bias, but ignores a shock that
costs $9$ when it occurs, with probability $q$. For which $q$ is the rule better?
\end{exercise}

\begin{exercise}[$\star\star\star$]\label{exo:b1:stabilisation-policy:refute}
What evidence would refute the claim that \omterm{def:b1:is-lm:fiscal-multiplier}{fiscal multipliers} are larger at
the \omterm{def:b1:is-lm:liquidity-trap}{zero lower bound}?
\end{exercise}

\section{Problem: Brenland's Finance Minister}\label{sec:b1:stabilisation-policy:problem}

\begin{problem}[{Weekend problem --- a forecast, a lag and the stabilisers}]\label{pb:b1:stabilisation-policy:brenland}
Brenland's output gap without policy follows $y_t = \rho y_{t-1} + e_t$, in per cent and by
quarter. A shock $e_1 = -5$ hits in quarter 1, with no shock afterwards. Policy, if
used, adds $p_t = -0.5y_{t-k}$ to the gap of quarter $t$: half the gap that would have
occurred $k$ quarters before without policy. The effect of each $p_t$ lasts one
quarter. The loss is the sum of squared gaps over quarters 1 to 8.

\medskip\noindent\textbf{Part I --- A forecast.} Take $\rho = 0.6$.
\begin{enumerate}
  \item Give the gap path for quarters 1 to 8 without policy.
  \item Compute the loss.
  \item The forecast made before quarter 1 was a shock of $-2$. Give the forecast path.
  \item What is the forecast error in quarter 1?
  \item How often did forecasts made in April of the year before foresee recessions in 1992--2014?
\end{enumerate}

\medskip\noindent\textbf{Part II --- Acting with a lag of two quarters.} Take $k = 2$.
\begin{enumerate}[resume]
  \item With $\rho = 0.6$, give the gap path with policy.
  \item Compute the loss, and compare with question 2.
  \item With $\rho = 0.9$, give the losses with and without policy.
  \item With $\rho = 0$, a shock that vanishes at once, give the losses.
  \item Explain the case $\rho = 0$.
  \item What must a minister know to decide whether to act?
\end{enumerate}

\medskip\noindent\textbf{Part III --- The automatic stabilisers.} Brenland's 2025
income tax has bands of $0\%$ to 12,000 dalens, $20\%$ to 40,000, $35\%$ to 100,000 and $45\%$
above; households with gross income below 12,000 receive $6{,}000 + 1{,}500(n - 1)$ dalens,
and every household $1{,}200$ per member beyond two. The survey has 2,000
households, 1,439 of them working.
\begin{enumerate}[resume]
  \item What tax does a household with 30,000 dalens pay?
  \item What does it pay if its income falls by $10\%$, and what is the
    fall in its \omterm{def:b1:national-accounts-in-depth:gni}{disposable income}?
  \item What is its marginal tax rate, and what share of its loss do taxes absorb?
  \item For a working household that falls below 12,000, why can \omterm{def:b1:national-accounts-in-depth:gni}{disposable income} fall by less still?
  \item Every working household loses $10\%$ of its earnings. The average fall in
    market income is 2,226.9 dalens per household; that of \omterm{def:b1:national-accounts-in-depth:gni}{disposable income}
    1,532.3. What share do the stabilisers offset?
  \item How much of it comes from taxes and how much from transfers?
  \item With a propensity to consume of $0.8$ out of \omterm{def:b1:national-accounts-in-depth:gni}{disposable income}, compare
    the fall in consumption with and without the stabilisers.
  \item Why does this offset need no forecast?
\end{enumerate}

\medskip\noindent\textbf{Part IV --- A rule.} Spending is tied to last quarter's gap: $k = 1$.
\begin{enumerate}[resume]
  \item With $\rho = 0.6$, give the gap path and the loss.
  \item With $\rho = 0.9$, give the loss.
  \item With $\rho = 0$, give the loss.
  \item Compare the rule's losses with those of Part II.
  \item Why can a rule act after one quarter where a decision takes two?
  \item What does the case $\rho = 0$ say about any policy that responds to past gaps?
  \item Which instrument of Part III responds within the quarter?
  \item State the result: what share of a fall in market income do Brenland's automatic stabilisers offset?
\end{enumerate}
\end{problem}
